\documentclass{report}
\usepackage{amsmath,amssymb}
\setlength{\parindent}{0mm}
\setlength{\parskip}{1em}
\begin{document}
\begin{center}
\rule{6in}{1pt} \
{\large Daniil Svyatskiy \\
{\bf New preconditioning strategy for Jacobian-Free solvers coupled with advanced discretization methods}}

MS B284 \\ Theoretical division \\ Los Alamos National Laboratory \\ Los Alamos \\ NM,87545
\\
{\tt dasvyat@lanl.gov}\\
Konstantin Lipnikov\\
David Moulton\end{center}

Numerical algorithms for nonlinear PDEs employ often implicit or explicit-implicit
time integration schemes that result in a system of nonlinear discrete
equations. Another key component of these algorithms is
the spatial discretization of PDEs. The current practice of using
unstructured meshes that conform to the physical domain or solution
features reduces accuracy
of the conventional discretization methods. New advanced discretization
methods that remain
accurate on unstructured meshes become more complex. However, the potential impact of
modern discretization methods on multi-physics simulations
has not been yet realized because the additional complexity
cannot be managed by existing nonlinear solvers.

Variably saturated flow in porous media is a critical process
in a wide variety of environmental applications, including
contaminant transport in the subsurface and carbon cycling in degrading
permafrost. This flow process is often modeled by the nonlinear
PDE which is referred to as the Richards' equation.
In Richard's formulation the flow model is
considered under the assumption that the gas phase is immobile.
This simplification also makes
Richards' equation a strongly nonlinear
parabolic PDE, creating significant
challenges for the performance of nonlinear solvers.

The implicit discretization in time of its mixed-form is the standard
approach for this type of applications since it enables large time steps
and locally conserve mass. This scheme
leads to a nonlinear discrete system of equations that must be solved
at each time step.
Typically, this solution is obtained by
discretizing PDEs and then solving the resulting system of
nonlinear discrete equations with a Newton-Raphson-type method.
This approach is widely used in the case of cartesian
grids with isotropic or grid-aligned anisotropic
permeabilities, for which the two-point flux approximation of
the flux is accurate.
The Jacobian matrix is computed by
differentiating the discrete nonlinear system of equations analytically
with respect to the discrete pressure unknowns. This discrete Jacobian
matrix is non-symmetric but positive definite, and for certain
upwinding formulations, naturally picks up the appropriate upwinding
needed for stability of the first-order derivatives.
However, complexity and cost do increase if complex models of porosity or
density are incorporated, particularly if these models require
interpolation of tabular data.
The geometric complexity of the subsurface environment
requires to handle non-orthogonal and unstructured
meshes. To discretize the flow model in these settings
advanced discretization methods are essential.
These methods offer accurate schemes on general meshes, but
close-from formulas for the analytic
Jacobian may not exist, e.g slope-limiting
methods, nonlinear discretization schemes. For these discrete systems
an implementation of the standard Newton-Raphson-type
nonlinear solvers becomes very problematic.
Although, this complexity in the discretization is manageable, it is problematic
for the direct analytic differentiating of the nonlinear discrete
system, as it is costly or overly complex to evaluate and the
resulting Jacobian matrix may be singular.
Therefore, development of Jacobian-Free solvers has attracted a
lot of interest in the past decade.


Several Jacobian-Free methods have become quite popular in the recent years.
Methods, such as the Jacobian-Free Newton-Krylov method, only require the
action of the Jacobian matrix on a vector, which is approximated
via the numerical computation of the Gateaux derivative.
Accelerated fixed-point methods such as Anderson
Mixing and Nonlinear Krylov Acceleration (NKA) methods are design to extend
the power of linear iterative schemes to the nonlinear case. In all
of these methods the preconditioner plays a crucial role for the
efficacy of a nonlinear solver. The standard preconditioning approach
is based on a linearized counterpart of the original discrete system.
The strong nonlinearity in the coefficients limits the
efficacy of that approach and results in the growth of nonlinear iterations.


We propose and analyze a new preconditioning strategy that is based on a
stable discretization
of the continuum Jacobian. The underlying idea in this work is to reverse
the order of the
discretization and linearization steps in the development of the
nonlinear solver. Specifically, by performing the linearization step
first, the analytic Jacobian of the continuum model is derived and
then analyzed to establish requirements for accurate and stable
discretizations.
This strategy allows us to take full
advantage of advanced discretization methods and control properties of
intermediate solutions.
For example, we may use monotone schemes for the diffusive term in the Jacobian and
upwinded schemes for the advective term. This not only leads to an efficient
preconditioner but also allows us to control its numerical properties.

The proposed strategy was tested in simulations of water infiltration into a partially
saturated layered medium in
steady-state and transient regimes on
non-orthogonal subsurface topology with
heterogeneous permeabilities and water retention models.
These experiments demonstrate that the new preconditioner improves
convergence of the existing
Jacobian-Free solvers 3-20 times.


\end{document}
